\section{埋めてはならない穴}
\label{sec:holes}

次の関数は、このライブラリの公開 API に存在しない。呼び出し側が研究のために実験するなら、別リポジトリの作業枝で行い、本書のモジュール名を使わない。

\begin{enumerate}
\item 多体の \(J\) を空間平均し、その平均をアインシュタイン方程式のソースとして返す関数。
\item 共有空間格子の上でマクスウェル方程式を時間発展させる関数。ファラデー係数を六つの生成子へ足す読みは、別の記録（\texttt{EMControl}）に属する。それを格子の PDE と同一視しない。
\item 層の半径から \(\alpha_s\) を返す関数。強い力の結合は導出されていない。
\item 粒子の生成、消滅、核融合、ベータ崩壊を一段の中で実行する関数。
\item チャート上の近日点余剰を、内在ローターの積分結果として返す関数。
\item \(G\)、\(\ell\)、\(m\)、\(a_0\)、微細構造定数を、許容錐または \(\gamma_s^{\mathrm{eq}}\) の零点から返す関数。
\item \(r\to 0\) の接触、合体、硬い反射。
\item 書かれた作用の \(E\) と、チャート上の力学的エネルギーを加算して「全エネルギー」と呼ぶ関数。
\item 量子重力稿のディラック方程式 \(\gamma^\mu\partial_\mu\Psi=-m\Psi\) を、多体状態の時間積分に使う関数。スピノルの記録は内在角の閉形式の代替ではない。
\end{enumerate}

穴が埋まったように見える数値は、第\ref{sec:contract}章のログが「作業」と書いた段の出力である。その出力を、Lean の定理番号の隣に置かない。
